% ============================================================
\chapter{Evaluación}
\label{cap:evaluacion}
% ============================================================

Este capítulo evalúa el sistema kubescan completo. Se parte de una
línea de base tomada del estado del arte, frente a la cual se contrastan los
resultados de cada capa; a continuación se presentan dichos resultados en el
mismo orden que las fases de entrenamiento descritas en el
Capítulo~\ref{cap:herramienta}. El desglose por clases del clasificador de grafo
sobre el conjunto de test (Sección~\ref{sec:gnn_por_clase}) caracteriza la
dirección de los errores residuales, y el contraste con una herramienta de
análisis estático externa (Sección~\ref{sec:validacion_externa}) delimita hasta
qué punto el corpus admite verificación independiente. Sobre esa base se
sintetiza una métrica de proceso completo que resume el comportamiento extremo a
extremo del sistema y se verifica el cumplimiento de los requisitos funcionales y
no funcionales definidos en el Capítulo~\ref{cap:requisitos}. El capítulo cierra
con la evaluación de usabilidad y aplicabilidad realizada con usuarios expertos.
Las métricas empleadas se justifican en detalle en la
Sección~\ref{sec:metricas}.

% -----------------------------------------------------------
\section{Línea de base del estado del arte}
\label{sec:linea_base}
% -----------------------------------------------------------

Antes de presentar los resultados propios, esta sección fija la línea de base
frente a la que se contrastan: los trabajos revisados en el
Capítulo~\ref{cap:estado} que abordan un problema comparable al de cada capa. La
comparación se organiza por capa, dado que la Capa~1 dispone de referencias
directas en la literatura de clasificación de \textit{IaC}, mientras que la
Capa~2 carece de trabajos previos estrictamente comparables. En ambos casos se
prioriza la comparación con resultados obtenidos sobre conjuntos de datos de
tamaño y naturaleza similares; los resultados propios que superan o contrastan
con esta línea de base se presentan en las secciones siguientes.

\subsection{Capa 1 en contexto}

Rahman et~al.\ \citep{Rahman2022}, cuyos datos forman la base del presente
trabajo, no publican modelos de clasificación supervisada sobre el mismo
conjunto: el objetivo de su trabajo era la construcción del conjunto de datos y el
análisis de la densidad de configuraciones erróneas, no la predicción, por lo que no
ofrece una cifra de referencia directa. Tampoco la literatura de \acs{IDS} sobre
UNSW-NB15 \citep{Moustafa2015} proporciona una: ese trabajo presenta el conjunto
de datos y compara exactitud y tasa de falsas alarmas de otros algoritmos, sin
consignar F1 macro ni resultados de Random Forest
(Sección~\ref{sec:ml_vulnerabilidades}). En consecuencia, el umbral
F1\,>\,0,85 de OE2 es un criterio de diseño propio y no una cota tomada de la
literatura, y la comparación informativa para la Capa~1 no es contra ese umbral
sino contra una línea de base trivial calculada sobre el propio corpus: la mejor
regla de una sola característica alcanza F1 macro\,=\,0,8502, frente al 0,9946
del modelo (Sección~\ref{sec:layer1}). Esa separación es consistente
con la naturaleza del problema: las configuraciones erróneas en manifiestos
\acs{YAML} son mayoritariamente discretas y acumulativas (más banderas activas
implica mayor riesgo), lo que hace que el espacio de características sea
relativamente separable una vez que se dispone de un conjunto de datos
suficientemente representativo.

\subsection{Capa 2 en contexto}

Según la revisión bibliográfica no sistemática realizada en este trabajo
(Capítulo~\ref{cap:estado}), no se ha localizado ningún trabajo previo que
aborde la clasificación supervisada de cadenas de ataque en clústeres
Kubernetes mediante \acs{GNN}, por lo que la Capa~2 no dispone de una
línea de base directamente comparable. Las referencias más cercanas son los
trabajos de \acs{GNN} para \ac{IDS} en red, como E-GraphSAGE
\citep{Lo2022}, que documenta tasas de detección superiores a las de
clasificadores tabulares clásicos sobre las re-publicaciones en formato NetFlow
de los conjuntos de referencia del área; y la literatura que documenta la
escasez de conjuntos de datos de grafos etiquetados en ciberseguridad
\citep{GNNSurvey2024}. Cabe puntualizar que esta última no establece un régimen
de tamaño concreto ni prescribe un protocolo de evaluación, de modo que no puede
invocarse como respaldo de la varianza observada en este trabajo; la validación
cruzada estratificada se adopta aquí como decisión metodológica propia. La
ventaja del componente estructural
se mide sobre el conjunto de test genuinamente no visto: P@5 pasa de 0,40
(solo Random Forest) a 0,80 (configuraciones con \ac{GNN}) y P@1 de 0,00 a 1,00
(Tabla~\ref{tab:ensemble_ablation}; Sección~\ref{sec:layer2} y siguientes),
indicando que la información estructural del grafo es decisiva para identificar
cadenas de ataque reales. La varianza entre pliegues de la fase final
(\(\sigma\)\,=\,0,160 para P@5) no se contrasta aquí con ninguna cota externa,
al no disponerse de una referencia que la establezca para conjuntos de este
tamaño; en este trabajo tiene, no obstante, una causa estructural identificada: el
particionado por familias concentra las cadenas de validación de cada pliegue en
1, 1, 4, 5 y 5~familias independientes respectivamente
(Sección~\ref{sec:layer2}), de modo que la desviación estándar entre pliegues
mide en buena parte la dispersión de dificultad entre familias. El análisis
multi-semilla (Sección~\ref{sec:multisemilla}) matiza esa lectura: cuatro de los
cinco pliegues puntúan de forma idéntica con las tres semillas ensayadas, pero el
único que varía —el pliegue~1, correspondiente a la familia
\texttt{badpods\_*}— es también el que más contribuye a esa desviación. En ese
pliegue concreto, la dificultad de la familia y la sensibilidad a la semilla no
resultan separables con los datos disponibles, de modo que la atribución
estructural de la varianza es parcial y no total.

% -----------------------------------------------------------
\section{Resultados de la Capa 1: Random Forest}
\label{sec:layer1}
% -----------------------------------------------------------

Esta sección presenta los resultados de los dos modelos de la Capa~1 descritos en
el Capítulo~\ref{cap:herramienta}: el clasificador binario y el modelo de severidad. Se
presentan primero las métricas de clasificación binaria y la importancia de las
características, y a continuación los resultados del modelo de severidad. Ambos
modelos comparten el mismo conjunto de 28~características extraídas por
manifiesto, descrito en la Sección~\ref{sec:construccion_dataset}.

\subsection{Clasificación binaria}

La Tabla~\ref{tab:rf_results} resume los resultados del modelo binario mediante
cuatro métricas: F1 macro (media no ponderada de F1 entre las clases seguro y
misconfigured, que da igual peso a ambas independientemente de su frecuencia en el
conjunto de datos), exactitud (\textit{accuracy}, porcentaje simple de aciertos), \ac{AUC}-\ac{ROC}
(capacidad del modelo para separar las dos clases a cualquier umbral de decisión) y
\textit{out-of-bag score} (estimación interna de generalización que
Random Forest calcula sobre las muestras que cada árbol no vio durante su
entrenamiento, sin necesitar un conjunto de validación aparte). El clasificador
alcanza una F1 macro de 0,9946 en el conjunto de test, superando el objetivo de
diseño (OE2, F1 macro\,>\,0,85) en 14,5 puntos porcentuales.

Conviene situar esa cifra frente a una línea de base trivial, porque el umbral de
OE2 resulta poco exigente sobre este corpus: la mejor regla de una sola
característica (\texttt{NO\_RESO}\,=\,1) alcanza por sí sola F1 macro\,=\,0,8502,
prácticamente el umbral de 0,85. El margen relevante del modelo no es, por tanto,
el que separa 0,9946 de 0,85, sino el que lo separa de esa regla trivial:
14,4~puntos porcentuales sobre la mejor decisión de una única condición. Una
disyunción de todas las banderas binarias, en cambio, se desploma a
F1 macro\,=\,0,3042, al marcar como mal configurado casi cualquier manifiesto.

\begin{table}[htbp]
\centering
\caption{Resultados del Random Forest — clasificación binaria}
\label{tab:rf_results}
\begin{tabular}{@{}lrrr@{}}
\toprule
Métrica & Test & \acs{CV} (5-fold) & Objetivo \\
\midrule
F1 macro & \textbf{0,9946} & 0,9917 \(\pm\) 0,0053 & \(>0{,}85\) \checkmark \\
Exactitud (\textit{accuracy}) & 0,9947 & 0,9919 \(\pm\) 0,0052 & — \\
\ac{AUC}-\ac{ROC} & \textbf{0,9986} & 0,9984 \(\pm\) 0,0011 & — \\
\textit{Out-of-bag score} & 0,9920 & — & — \\
\bottomrule
\end{tabular}
\fuentepropia
\end{table}

La Figura~\ref{fig:rf_confusion} muestra la matriz de confusión sobre el conjunto
de test (566~muestras). El modelo comete solo 3~errores totales: 3~falsos positivos
(manifiestos seguros clasificados como misconfigured) y 0~falsos negativos.
La ausencia de falsos negativos es la propiedad más crítica para seguridad operativa:
significa que ningún manifiesto con configuraciones erróneas reales escapa al filtrado de
la Capa~1, independientemente de cuántas banderas estén activas o de su rareza en el
conjunto de datos. Los 3~falsos positivos corresponden a manifiestos seguros que acumulan varias
features de riesgo de severidad media (por ejemplo, \texttt{NO\_RUN\_AS\_NON\_ROOT}
y \texttt{IMAGE\_USES\_LATEST} simultáneamente), pero que no tienen ninguna
bandera de escape activa; el operador de seguridad puede descartarlos rápidamente
inspeccionando el campo \texttt{flags} de la salida \ac{JSON}.

El \ac{AUC}-\ac{ROC} de 0,9986 (prácticamente~1,00) indica que el modelo puede ordenar las
muestras por probabilidad de configuración errónea de forma casi perfecta a cualquier umbral
de decisión. En la práctica, este valor implica que, si se ordenan los 566~manifiestos
de test de mayor a menor \textit{risk\_score}, los manifiestos misconfigured aparecen
casi todos antes que los seguros. Esta propiedad es la que habilita el uso del
\textit{risk\_score} como característica nodal cuantitativa en el grafo de la Capa~2:
la señal transmitida a la \ac{GNN} no es un bit binario sino una puntuación
continua que ordena los manifiestos por densidad de configuraciones erróneas. No se
aplica ninguna técnica de calibración de probabilidades, por lo que el valor no
debe interpretarse como una probabilidad calibrada en sentido estricto; el
ordenamiento resultante es, en cambio, consistente con la matriz de confusión de
la Figura~\ref{fig:rf_confusion}.

\begin{figure}[htbp]
  \centering
  \caption{Matriz de confusión del Random Forest en el conjunto de test}
  \label{fig:rf_confusion}
  \includegraphics[width=0.52\textwidth]{fig_rf_confusion}

  \fuentepropia
\end{figure}

\subsection{Importancia de características}
\label{sec:rf_importance}

La Figura~\ref{fig:rf_importance} muestra las 10 características más relevantes del
modelo. La característica más discriminativa es \texttt{total\_misconfigs}
(33,05\,\% de importancia), seguida de \texttt{NO\_RESO} (14,79\,\%) e
\texttt{INSECURE\_HTTP} (11,90\,\%). Las banderas de escape críticas (\texttt{TRUE\_HOST\_PID},
\texttt{SEC\_CONT\_OVER\_PRIVIL}) presentan importancias menores (\(<\)1\,\%) por su
rareza en el conjunto de datos, pero son identificadas correctamente cuando están presentes,
como confirma la ausencia de falsos negativos en la clasificación binaria. Las
dos primeras características de la Figura~\ref{fig:rf_importance} (en rojo)
concentran el 48\,\% de la capacidad discriminativa total del modelo. Un estudio
de ablación eliminando \texttt{total\_misconfigs} produce F1 idéntico y la misma
matriz de confusión, lo que confirma que el modelo no depende de esa
característica derivada: su señal es recuperable a partir de las banderas que la
componen.

Esa ablación, sin embargo, no permite concluir que no exista circularidad entre
etiqueta y características. La etiqueta binaria del corpus de Rahman et~al.\ se
define como la presencia de al menos una configuración errónea de un subconjunto
concreto de patrones, y esos patrones son a su vez características de entrada del
modelo: sobre las 1.906~filas de esa fuente, la etiqueta coincide exactamente con
la disyunción de ocho de las características observadas
(\texttt{INSECURE\_HTTP}, \texttt{NO\_SECU\_CONTEXT},
\texttt{WITHIN\_MANIFEST\_SECRET}, \texttt{NO\_RESO}, \texttt{CAP\_SYS\_ADMIN},
\texttt{TRUE\_HOST\_NET}, \texttt{ALLOW\_PRIVI} y \texttt{DOCKERSOCK\_PATH}), sin
un solo falso positivo ni falso negativo. Reentrenar la misma configuración
restringida a esas filas produce, consecuentemente, F1 macro\,=\,1,0000. La Capa~1
aprende por tanto una aproximación robusta de una regla determinista definida
sobre su propio espacio de características —bajo imputación, valores ausentes y
las etiquetas heterogéneas del resto de fuentes—, no una noción de
configuración errónea independiente de dicha regla. Esta limitación, análoga a la que
afecta al etiquetado de grafos, se analiza en la limitación~L5 del
Capítulo~\ref{cap:conclusiones}.

\begin{figure}[htbp]
  \centering
  \caption{Importancia de características del clasificador Random Forest}
  \label{fig:rf_importance}
  \includegraphics[width=\textwidth]{fig_rf_importance}

  \fuentepropia
\end{figure}

\subsection{Modelo de severidad}

El clasificador de tres clases (\textit{clean}, \textit{low-medium},
\textit{high-critical}) alcanza F1 macro\,=\,0,9838 en test. La Tabla~\ref{tab:severity_f1}
desglosa por clase la precisión (proporción de predicciones de esa clase que son
correctas), el \textit{recall} (proporción de casos reales de esa clase que el
modelo identifica), el F1 (media armónica de ambas) y el soporte (número de
muestras de test en esa clase):

\begin{table}[htbp]
\centering
\caption{F1 por clase — modelo de severidad (3 clases, test)}
\label{tab:severity_f1}
\begin{tabular}{@{}lrrrr@{}}
\toprule
Clase & Precisión & Recall & F1 & Soporte \\
\midrule
\textit{clean}          & 1,0000 & 0,9906 & 0,9953 & 319 \\
\textit{low-medium}     & 0,9615 & 0,9934 & 0,9772 & 151 \\
\textit{high-critical}  & 0,9894 & 0,9688 & 0,9789 & 96 \\
\midrule
Macro media & 0,9836 & 0,9842 & 0,9838 & — \\
\bottomrule
\end{tabular}
\fuentepropia
\end{table}

El F1 más bajo de las tres clases corresponde a \textit{low-medium} (0,9772), por
debajo del de \textit{high-critical} (0,9789). La causa no es una dificultad
intrínseca de la clase intermedia, cuyo \textit{recall} es de 0,9934, sino su
papel de sumidero de los errores del modelo: absorbe los seis manifiestos mal
asignados del conjunto de test —tres procedentes de \textit{clean} y tres de
\textit{high-critical}—, lo que rebaja su precisión a 0,9615. La clase
\textit{high-critical} (96~muestras de test, que representan manifiestos con
banderas de escape activas como \texttt{TRUE\_HOST\_PID} o
\texttt{SEC\_CONT\_OVER\_PRIVIL}) comete 3~errores, los tres clasificados como
\textit{low-medium}; ningún manifiesto crítico es clasificado como
\textit{clean}. Esa ausencia de fugas hacia la clase \textit{clean} es la
propiedad relevante desde el punto de vista de seguridad operativa: en el peor
caso un manifiesto crítico se subestima a \textit{low-medium}, nunca se descarta
como seguro, y la dirección de los errores restantes —tres manifiestos
\textit{clean} sobreclasificados como \textit{low-medium} y un único manifiesto
\textit{low-medium} elevado a \textit{high-critical}— es igualmente conservadora.
El modelo de severidad se utiliza en el informe de salida
de kubescan para proporcionar al operador una descripción cualitativa
del nivel de riesgo de cada manifiesto, complementando el \textit{risk\_score}
numérico del modelo binario.

% -----------------------------------------------------------
\section{Resultados de la Capa 2: \acs{GNN}}
\label{sec:layer2}
% -----------------------------------------------------------

Esta sección presenta los resultados de la Capa~2 en el orden en que fueron
obtenidos durante el desarrollo experimental: la evolución de las métricas a lo
largo de las fases de mejora del modelo, el análisis de Precisión@5, los
resultados por pliegue de la configuración final, la ablación de la función de
pérdida y del criterio de selección, y la robustez frente a la semilla de
entrenamiento. A diferencia de la Capa~1, la Capa~2 requirió varias iteraciones
metodológicas antes de alcanzar una configuración cuyas métricas fueran a la vez
válidas y reproducibles. Cada iteración corrigió una fuente concreta de fuga de
información o una elección de entrenamiento subóptima, y el descenso de las
métricas que acompaña a esas correcciones cuantifica el sesgo optimista que
introducían los protocolos de particionado sucesivamente descartados.

\subsection{Evolución por fase experimental}

La Tabla~\ref{tab:gnn_evolution} muestra la evolución de los resultados de la
\ac{GAT} a lo largo de las seis fases del desarrollo experimental. En las cuatro
primeras, la aumentación de datos y la corrección del vector de escape
\texttt{HOSTPATH\_MOUNT} contribuyeron de forma independiente y complementaria. La quinta fase corrige el protocolo de particionado: las fases~1--4
incluían en el entrenamiento variantes aumentadas derivadas de los
clústeres de validación y test (fuga de datos por casi-duplicados), por lo que sus
métricas están infladas de forma optimista. Esa quinta fase aplica el protocolo
\textit{group-aware} descrito en el Capítulo~\ref{cap:herramienta} (las variantes acompañan a su
clúster base y los clústeres de test quedan fuera de los pliegues) e incorpora el
\textit{embedding} de tipos de arista; sus métricas son metodológicamente válidas
pero corresponden todavía al corpus preliminar de 471~grafos. La sexta y última
fase incorpora el corpus completo de Rahman et~al.\ (599~grafos originales,
37~cadenas; Sección~\ref{sec:construccion_dataset}) y extiende el protocolo de
particionado a las \textbf{familias de plantilla} (Capítulo~\ref{cap:herramienta}):
durante esta fase se detectó que los casi duplicados de una misma plantilla
(ocho \texttt{badpods\_*}, nueve \texttt{datadog\_*}, ocho
\texttt{simulator\_*}, doce \texttt{kubernetes\_goat\_*}) quedaban repartidos
entre particiones, una fuga de información adicional a la de las variantes
aumentadas que inflaba las métricas de las fases 1--5. Con las familias atómicas,
entropía cruzada ponderada como función de pérdida y selección de
\textit{checkpoints} por Precisión@5 (ablación en la
Sección~\ref{sec:ablacion_perdida}), esta fase es la que se presenta como
resultado final del trabajo: sus métricas son las que produce el modelo
efectivamente desplegado y reproducible a partir de los \textit{checkpoints} del
repositorio. Por el cambio de protocolo, sus cifras no son directamente
comparables con las de las fases~1--5 de la Tabla~\ref{tab:gnn_evolution}: miden
generalización entre familias, no interpolación entre casi duplicados.

\begin{table}[htbp]
\centering\small
\caption{Evolución de resultados de la Capa~2 (\acs{GAT}, 5-fold \acs{CV})}
\label{tab:gnn_evolution}
\begin{tabular}{@{}p{4.3cm}rrrrc@{}}
\toprule
Fase & Grafos & Cadenas & F1 macro & P@5 & \shortstack{Cobert.\\techo} \\
\midrule
Línea base & 82 & 13 & 0,829 \(\pm\) 0,098 & 0,400 \(\pm\) 0,179 & 2/5 \\
Tras aumentación & 277 & 13 & 0,915 \(\pm\) 0,059 & 0,520 \(\pm\) 0,098 & 5/5 \\
Tras corrección HOSTPATH & 307 & 15 & 0,915 \(\pm\) 0,050 & 0,600 \(\pm\) 0,000 & 5/5 \\
Conjunto de datos extendido & 471 & 25 & 0,917 \(\pm\) 0,065 & 0,880 \(\pm\) 0,098 & 5/5 \\
Protocolo \textit{group-aware} + \textit{embedding} de aristas
  & 471 & 25 & 0,870 \(\pm\) 0,075 & 0,720 \(\pm\) 0,098 & 3/5 \\
\addlinespace
Corpus completo + familias de plantilla (final)
  & 599 & 37 & \textbf{0,803 \(\pm\) 0,137} & \textbf{0,720 \(\pm\) 0,160} & 1/5 \\
\bottomrule
\end{tabular}
\fuentepropia
\end{table}

En la Tabla~\ref{tab:gnn_evolution}, las fases 1--4 emplean el protocolo de
particionado preliminar (con fuga por variantes aumentadas) y se conservan como
registro histórico; la fase~5 introduce el protocolo \textit{group-aware} sobre el
corpus preliminar; la fase~6 es el resultado final, con protocolo
\textit{group-aware} sobre el corpus completo.
La columna «Cobert.\ techo» indica en cuántos de los 5~pliegues el modelo alcanza
el \emph{techo teórico} de P@5. La métrica se calcula como el número de cadenas
situadas en el top-5 dividido entre \(\min(5,\, n_{\text{pos}})\), siendo
\(n_{\text{pos}}\) el número de cadenas presentes en el pliegue, de modo que el
techo vale 1,0 siempre que el pliegue contenga al menos una cadena y se alcanza
cuando las cinco primeras posiciones del \textit{ranking} están ocupadas por
cadenas —o, si el pliegue contiene menos de cinco, cuando todas ellas figuran en
el top-5. En la fase final, los cinco pliegues de validación contienen 8, 7, 6, 6
y 5~cadenas respectivamente, por lo que el denominador es 5 en todos ellos y el
techo 1,0 exige cinco aciertos en las cinco primeras posiciones; únicamente el
pliegue~0 lo consigue.

Las cuatro primeras fases de la Tabla~\ref{tab:gnn_evolution} corresponden al
protocolo preliminar. El objetivo de diseño de \textit{ranking} es P@5\,>\,0,70 y el umbral
de F1 fijado para la Capa~1, tomado como referencia, 0,85; las desviaciones típicas
representan la desviación estándar entre los 5~pliegues de validación cruzada.

\subsection{Análisis de Precisión@5}

El objetivo de diseño es P@5\,>\,0,70 (OE4). En el resultado final de este trabajo
—particionado consciente de familias sobre el corpus completo de Rahman et~al.
(599~grafos originales, 37~cadenas; fase~6 de la Tabla~\ref{tab:gnn_evolution})—
los pliegues de validación cruzada particionan los 513~grafos originales no
reservados para test (32~cadenas). El modelo obtiene
P@5\,=\,0,720\,\(\pm\)\,0,160, alcanzando el objetivo con la semilla de
entrenamiento desplegada (42); en las tres semillas (7, 42, 123) la media es
de 0,68\,\(\pm\)\,0,03\footnote{Las desviaciones típicas de las métricas del modelo
se expresan como desviación poblacional (\(\text{ddof}=0\)), consistente con la
instrumentación de NumPy empleada en la validación cruzada; la del \ac{SUS}
(Sección~\ref{sec:usabilidad}) sigue la convención muestral (\(n-1\)) habitual en su
literatura.}, con el objetivo dentro de una desviación típica de la
media (Sección~\ref{sec:multisemilla}).

Alcanzar esta cifra exigió diagnosticar una regresión aparente: la primera
configuración entrenada sobre el corpus completo con familias atómicas empleaba
\textit{focal loss} y selección de \textit{checkpoints} por F1 macro, y obtuvo
P@5\,=\,0,480\,\(\pm\)\,0,271, muy por debajo tanto del objetivo como de la
fase~5. La ablación factorial de la Sección~\ref{sec:ablacion_perdida} mostró
que la caída se debía a las elecciones de función de pérdida y de criterio de
selección —no al corpus ni al protocolo de particionado—, y que revertir a
entropía cruzada ponderada y seleccionar por Precisión@5 recupera el objetivo
sobre el mismo particionado honesto. La varianza entre pliegues
(\(\sigma\)\,=\,0,160, frente a 0,098 en la fase~5) refleja en gran medida la
propia atomicidad de las familias: cada pliegue evalúa 1, 1, 4, 5 y
5~familias independientes respectivamente, como se detalla a continuación.

\subsection{Resultados por pliegue (fase final)}

Las medias de validación cruzada de la configuración desplegada ocultan una
dispersión considerable entre pliegues, y su desglose es necesario para atribuir
esa dispersión a una causa concreta. La Tabla~\ref{tab:gnn_folds} recoge la F1
macro y la Precisión@5 obtenidas en cada uno de los cinco pliegues, junto con la
media y la desviación típica poblacional que resumen el conjunto. Los valores
individuales se interpretan a continuación a la luz de la composición de cada
pliegue en términos de familias de plantilla.

\begin{table}[htbp]
\centering
\caption{Resultados por pliegue — Capa~2, fase final}
\label{tab:gnn_folds}
\begin{tabular}{@{}crr@{}}
\toprule
Pliegue & F1 macro & P@5 \\
\midrule
0 & 1,000 & 1,00 \\
1 & 0,580 & 0,60 \\
2 & 0,863 & 0,60 \\
3 & 0,801 & 0,60 \\
4 & 0,770 & 0,80 \\
\midrule
\textbf{Media} & \textbf{0,803} & \textbf{0,720} \\
\textbf{\(\pm\) Desv.} & \textbf{0,137} & \textbf{0,160} \\
\bottomrule
\end{tabular}
\fuentepropia
\end{table}

La Tabla~\ref{tab:gnn_folds} corresponde al particionado consciente de familias
sobre el corpus completo: los pliegues particionan 513~grafos originales con
32~cadenas, con el conjunto de test excluido.
La dispersión entre pliegues es consecuencia directa del particionado por
familias: al mantener cada familia de plantilla como unidad atómica, las cadenas
de validación de cada pliegue quedan concentradas en pocas familias
independientes. El pliegue~0 valida sobre las nueve variantes \texttt{datadog\_*}
(ocho de ellas cadenas; grafos de 2--3 nodos casi idénticos entre sí) y alcanza
F1\,=\,1,0 y P@5\,=\,1,0; el pliegue~1 valida sobre las ocho variantes
\texttt{badpods\_*} (siete de ellas cadenas) y
es sistemáticamente el más difícil (P@5\,=\,0,60 con la semilla desplegada;
0,2--0,6 según la semilla, Sección~\ref{sec:multisemilla}). Los tres pliegues
restantes reparten las dos familias grandes que quedan junto a repositorios
únicos: el pliegue~2 concentra las ocho variantes \texttt{simulator\_*}, el
pliegue~3 las doce variantes \texttt{kubernetes\_goat\_*} y el pliegue~4 no
contiene ninguna de las cuatro familias de plantilla, sino exclusivamente
repositorios únicos. El número de familias independientes portadoras de cadena
por pliegue es, en consecuencia, de 1, 1, 4, 5 y 5. Cada pliegue mide, por tanto,
la capacidad del modelo de generalizar a un número reducido de familias no vistas
—un tamaño muestral efectivo menor que el número nominal de cadenas—, lo que
explica que la varianza
(\(\sigma\)\,=\,0,160) sea mayor que bajo protocolos con fuga y constituye una
amenaza a la validez de conclusión que se documenta en el
Capítulo~\ref{cap:conclusiones}.

\subsection{Ablación de función de pérdida y criterio de selección}
\label{sec:ablacion_perdida}

La regresión aparente de la primera configuración sobre el corpus completo
(P@5\,=\,0,480) confundía dos cambios simultáneos: la introducción de
\textit{focal loss} \citep{Lin2017} y el nuevo particionado por familias. Para
atribuir el efecto se completó la matriz factorial 2\(\times\)2 —función de
pérdida (entropía cruzada ponderada frente a \textit{focal loss},
\(\gamma\)\,=\,2) por criterio de selección de \textit{checkpoints} (F1 macro
frente a Precisión@5)— con particionado, semilla y arquitectura idénticos
(Tabla~\ref{tab:loss_ablation}).\footnote{Las cuatro configuraciones son reproducibles
mediante \texttt{make reproduce GNN\_FOCAL=\{0,1\} GNN\_SELECT\_BY=\{p5,f1\}}; solo se
conserva el \textit{checkpoint} de la configuración finalmente desplegada.}

\begin{table}[htbp]
\centering
\caption{Ablación de función de pérdida y criterio de selección}
\label{tab:loss_ablation}
\begin{tabular}{@{}llrr@{}}
\toprule
Pérdida & Selección & P@5 (\acs{CV}) & F1 macro (\acs{CV}) \\
\midrule
\textit{Focal} (\(\gamma\)=2) & F1 macro & 0,480 \(\pm\) 0,271 & 0,823 \(\pm\) 0,097 \\
Entropía cruzada pond. & F1 macro & 0,600 \(\pm\) 0,283 & 0,859 \(\pm\) 0,103 \\
\textit{Focal} (\(\gamma\)=2) & P@5 & 0,680 \(\pm\) 0,204 & 0,795 \(\pm\) 0,173 \\
Entropía cruzada pond. & P@5 & \textbf{0,720 \(\pm\) 0,160} & 0,803 \(\pm\) 0,137 \\
\bottomrule
\end{tabular}
\fuentepropia
\end{table}

Las cuatro configuraciones de la Tabla~\ref{tab:loss_ablation} se evaluaron con
5-fold \ac{CV} sobre el corpus completo, particionado por familias y semilla~42.
Los dos efectos principales apuntan en la misma dirección, pero no se comportan
de forma aditiva: sustituir \textit{focal loss} por entropía cruzada ponderada
aporta +0,12~puntos de P@5 con selección por F1 macro y solo +0,04 con selección
por P@5, mientras que seleccionar los \textit{checkpoints} por la métrica
primaria de \textit{ranking} en lugar de por F1 macro aporta +0,20~puntos bajo
\textit{focal loss} y +0,12 bajo entropía cruzada. La suma de ambos efectos
principales sobre la configuración de partida (0,480) predeciría
P@5\,=\,0,80, frente al 0,720 efectivamente observado: existe una interacción
negativa de 0,08~puntos, es decir, la mejora que aporta cada factor se solapa
parcialmente con la del otro en lugar de acumularse. La \textit{focal loss},
diseñada para regímenes con miles de ejemplos por clase \citep{Lin2017}, además
anula por completo la clasificación de la clase \textit{attack\_chain} en el
conjunto de test (F1 de clase 0,0 frente a 0,667 con entropía cruzada). La
configuración desplegada es, en consecuencia, entropía cruzada ponderada con
selección por Precisión@5.

La Tabla~\ref{tab:loss_ablation} recoge la Precisión@5 de validación cruzada
(\(\pm\)\(\sigma\)) de las cuatro configuraciones; solo la desplegada —entropía
cruzada ponderada con selección por Precisión@5— supera el objetivo de diseño.
La distancia respecto de la segunda mejor configuración es, no obstante, estrecha
(0,720\,\(\pm\)\,0,160 frente a 0,680\,\(\pm\)\,0,204): sus bandas de una
desviación típica se solapan casi por completo y, sobre cinco pliegues y con una
métrica cuantizada en múltiplos de 0,2, esa diferencia equivale a un solo acierto
de \textit{ranking} en un único pliegue. Los datos sostienen, por tanto, una
afirmación más débil que la de una dependencia estricta de ambos factores: la
combinación desplegada es la única que sitúa la media de validación cruzada por
encima del objetivo, y la configuración más alejada de él —\textit{focal loss}
con selección por F1 macro, 0,480— sí queda separada por un margen que excede la
dispersión observada.

\subsection{Ablación de arquitectura}
\label{sec:ablacion_arquitectura}

La elección de \ac{GAT} frente a alternativas más simples y la profundidad de
tres capas se fijaron en el diseño de la Capa~2 y admiten contraste empírico
sobre el mismo protocolo. La literatura sobre \acs{GNN} aplicadas a grafos
pequeños heterogéneos documenta de forma consistente que los mecanismos de
atención mejoran la precisión cuando los tipos de arista tienen importancias
desiguales \citep{Velickovic2018}. En el grafo de clúster
Kubernetes, las aristas de \textit{privilege\_reach} (tipo~1) son
cualitativamente más informativas que las de \textit{dir\_proximity} (tipo~0): la
primera conecta directamente nodos de escape con el resto del clúster, mientras que
la segunda simplemente refleja la organización de directorios del repositorio.
\ac{GAT} puede aprender esta distinción durante el entrenamiento, mientras que \ac{GCN} asigna
pesos uniformes a todos los vecinos, diluyendo la señal de escape con ruido de
co-localización.

Esta hipótesis se cuantifica empíricamente mediante un experimento de ablación
controlado sobre el corpus completo, con particionado, semilla, función de
pérdida y criterio de selección idénticos a los de la configuración final
(fase~6 de la Tabla~\ref{tab:gnn_evolution}), de modo que las tres filas son
directamente comparables entre sí y con el resultado final (Tabla~\ref{tab:arch_ablation}):

\begin{table}[htbp]
\centering
\caption{Ablación de arquitectura — Capa~2}
\label{tab:arch_ablation}
\begin{tabular}{@{}lrr@{}}
\toprule
Arquitectura & F1 macro & P@5 \\
\midrule
\ac{GAT} 3 capas + \textit{embedding} de aristas (desplegada)
  & \textbf{0,803 \(\pm\) 0,137} & \textbf{0,720 \(\pm\) 0,160} \\
\ac{GAT} 2 capas + \textit{embedding} de aristas
  & 0,800 \(\pm\) 0,171 & \textbf{0,720 \(\pm\) 0,160} \\
\ac{GCN} 3 capas (sin tipos de arista)
  & 0,711 \(\pm\) 0,143 & 0,680 \(\pm\) 0,160 \\
\bottomrule
\end{tabular}
\fuentepropia
\end{table}

La \ac{GAT} supera a la \ac{GCN} en F1 (+9,2~puntos) y en P@5 (+4~puntos). Esa
ventaja acredita la configuración adoptada en su conjunto, pero no permite
atribuirla al esquema de tipos de arista por separado: la rama \ac{GCN} sustituye
simultáneamente el mecanismo de agregación por atención y descarta el tensor
\texttt{edge\_attr}, de modo que ambos factores varían a la vez. Aislar la
contribución del tipado exigiría una tercera rama —\ac{GAT} sin tipos de arista—
que no forma parte de las configuraciones entrenadas
(Capítulo~\ref{cap:conclusiones}). La comparación entre 2 y 3~capas es
menos concluyente: las medias de validación cruzada casi coinciden
—P@5\,=\,0,720\,\(\pm\)\,0,160 en ambas configuraciones y F1 macro de 0,803
frente a 0,800—, pero esa coincidencia agregada oculta un desacuerdo apreciable
pliegue a pliegue, con diferencias de F1 de \(+\)0,142 en el pliegue~3 y
\(-\)0,191 en el pliegue~4 que se cancelan al promediar y dejan una brecha media
de solo 0,003. Dos modelos que difieren en esa magnitud sobre pliegues concretos
no evidencian robustez de la arquitectura frente a la profundidad, sino que la
profundidad no determina el nivel medio de las métricas sobre este corpus. Se
mantiene la configuración de 3~capas como desplegada por su ventaja marginal en
F1 macro y por coherencia con el resto de resultados presentados. La
ablación de \textit{pooling} (\textit{mean}-only frente a \textit{mean+max}) queda
como trabajo futuro (TF6, Sección~\ref{sec:trabajo_futuro}).

\subsection{Robustez frente a la semilla de entrenamiento}
\label{sec:multisemilla}

Para acotar la dependencia del resultado respecto de la semilla de entrenamiento,
la configuración final se reentrenó con tres semillas (7, 42, 123;
Figura~\ref{fig:multiseed_folds}) manteniendo fijo el particionado (semilla~42):
re-sortear también las particiones confundiría la varianza del modelo con la
asignación de familias a pliegues, que la Sección~\ref{sec:layer2} muestra
dominante. Las métricas de test aplican la restricción de viabilidad estructural
(Sección~\ref{sec:gate_estructural}).

Dos observaciones acotan la lectura del 0,72. Primera: la variación entre
semillas se concentra íntegramente en el pliegue~1 (la familia
\texttt{badpods\_*}: P@5 de 0,2/0,4/0,6 según semilla), mientras que los otros
cuatro pliegues puntúan de forma idéntica con las tres semillas. Esa coincidencia
impide leer el experimento como una demostración de estabilidad global, porque el
único pliegue sensible a la semilla es precisamente el que domina la desviación
típica entre pliegues: en él, la dificultad de la familia y la variabilidad de
entrenamiento quedan confundidas y no pueden separarse con tres semillas. Lo que
el experimento sí establece es que cuatro quintas partes del particionado son
invariantes a la semilla y que la incertidumbre residual de la estimación de
validación cruzada se localiza en una sola familia de plantilla. Segunda: las
métricas de test a nivel de despliegue son estables —P@1\,=\,1,00 y
P@3\,=\,1,00 con las tres semillas, y P@5 entre 0,80 y 1,00—, de modo que la
sensibilidad a la semilla afecta a la estimación de validación cruzada, no al
comportamiento observable del sistema sobre repositorios reales no vistos.

\begin{figure}[htbp]
  \centering
  \caption{Precisión@5 por pliegue con tres semillas de entrenamiento}
  \label{fig:multiseed_folds}
  \includegraphics[width=0.85\textwidth]{fig_multiseed_folds}

  \fuentepropia
\end{figure}

% -----------------------------------------------------------
\section{Resultados de la Capa 3: \textit{Ensemble}}
\label{sec:layer3}
% -----------------------------------------------------------

El ajuste de los pesos del \textit{ensemble} atravesó dos correcciones metodológicas
sucesivas, ambas implementadas en \texttt{run\_ga\_ensemble.py}. En una primera
iteración, el \ac{GA} original convergía sistemáticamente a pesos casi uniformes
\((w_{\text{rf}}, w_{\text{gnn}}, w_{\text{esc}}) \approx (0{,}34,\, 0{,}33,\, 0{,}33)\)
sin mejorar sobre su población inicial en 150~generaciones, un patrón que parecía
indicar una meseta plana en la función de \textit{fitness}. Un diagnóstico
descartó esa hipótesis: el cruce por promedio aritmético es contractivo y no
puede alcanzar los vértices del símplex de pesos desde una población inicializada
por muestreo de Dirichlet, y la mutación gaussiana (\(\sigma=0{,}08\)) es
demasiado local para recorrer esa distancia en 150~generaciones. La primera
corrección sembró los vértices y aristas del símplex en la población inicial y
añadió una mutación de salto de frontera de baja probabilidad; con este cambio,
el \ac{GA} pasó a converger de forma determinista, en las 5~semillas
independientes probadas, al mismo punto que una búsqueda en rejilla
determinística (Sección~\ref{sec:capa3_diseno}, 231~evaluaciones): el vértice de
señal de escape pura (\(w_{\text{esc}}=1\), \textit{objective}\,=\,0,86),
confirmando que la función objetivo tenía un óptimo global único y alcanzable, no
una meseta degenerada. Sin embargo, ese óptimo \textit{out-of-fold} no
generalizaba: al puntuar únicamente por la señal binaria de escape, el \textit{ranking} se
colapsaba a un grupo de empates que fallaba estructuralmente ante cualquier
cadena de ataque cuya única señal fuerte no fuera la bandera de escape.

Esta segunda observación motivó la corrección definitiva: regularizar la propia
función objetivo con una penalización de piso,
\(F = \alpha \cdot P@5 + \beta \cdot (1-\text{\ac{FPR}\_clean}) -
\lambda \sum_i \max(0,\, w_{\min} - w_i)\), con \(w_{\min}=0{,}1\) y
\(\lambda=2{,}0\) (opciones \texttt{--min-weight}/\texttt{--reg-lambda}). Esta
penalización hace que el vértice de escape puro deje de ser competitivo frente a
soluciones interiores del símplex. Sobre el conjunto \textit{out-of-fold} de la
configuración final (513~grafos, 32~cadenas), el \ac{GA} regularizado encuentra
su mejor solución ya en la generación inicial y no la mejora durante las
150~generaciones:
\((w_{\text{rf}}, w_{\text{gnn}}, w_{\text{esc}}) = (0{,}517,\, 0{,}115,\, 0{,}368)\),
\textit{objective}\,=\,0,72. La búsqueda en rejilla encuentra un punto muy
distinto (\(w_{\text{rf}}=0{,}80,\, w_{\text{gnn}}=0{,}10,\,
w_{\text{esc}}=0{,}10\)) con idéntica P@5 \textit{out-of-fold}: al ser P@5 una
función escalonada (solo toma valores múltiplos de 0,2), la función objetivo
presenta mesetas extensas en las que los pesos quedan \textbf{indeterminados}.
Esta indeterminación es medible: sobre el conjunto de test, los pesos elegidos
por el \ac{GA} y los pesos uniformes \((1/3)\) producen métricas idénticas
(Tabla~\ref{tab:ensemble_ablation}). Contribuye a ella un desplazamiento de
distribución entre el conjunto de ajuste y el de test: 15 de las 32~cadenas
\textit{out-of-fold} pertenecen a familias de \textit{fixtures} sintéticas
(\texttt{badpods\_*}, \texttt{datadog\_*}) cuyos manifiestos reciben
\textit{risk scores} altos del Random Forest, mientras que las cinco cadenas de test son
repositorios reales cuyo \textit{risk score} medio por clúster es
\(\approx\)0 (la media se diluye entre manifiestos mayoritariamente benignos).
Tres refinamientos del ajuste (un término de desempate por \textit{mean
reciprocal rank}, la agregación del \textit{pool} por familias y el ajuste sobre
el \textit{ranking} con restricción estructural) se implementaron y evaluaron
sin que ninguno mejorara las métricas de test, por lo que los pesos se ajustan
sobre el \textit{ranking} sin restricción y la restricción de viabilidad
estructural (Sección~\ref{sec:gate_estructural}) se aplica en inferencia y
evaluación.

La evaluación final se realiza sobre el conjunto de test (86~grafos: 3~limpios,
78~\textit{isolated} y 5~cadenas de ataque; techo P@5\,=\,1,00). Bajo el
particionado consciente de familias, las cinco cadenas son repositorios reales
únicos —sin hermanos de plantilla en el entrenamiento— excluidos de los pliegues
de validación cruzada, del ajuste de pesos del \ac{GA} y del entrenamiento con
variantes aumentadas: no vistos en ningún momento del desarrollo. Con la
restricción de viabilidad estructural aplicada al \textit{ranking}
(Sección~\ref{sec:gate_estructural}), el \textit{ensemble} alcanza P@1\,=\,1,00,
P@3\,=\,1,00 y P@5\,=\,0,80: las cadenas ocupan las posiciones 1--4 y 6, y solo
\texttt{minik8s-ctf} queda fuera del top-5. Sin la restricción, las métricas
caen a P@1\,=\,0,00, P@3\,=\,0,33 y P@5\,=\,0,60 (Figura~\ref{fig:structural_gate}): la regla es decisiva porque
74 de los 86~grafos de test son de un solo nodo, y dos de ellos, con todas sus
banderas de escape activas, encabezarían el \textit{ranking} pese a no poder
albergar una cadena. Procede delimitar, en consecuencia, cuál es el conjunto sobre
el que se ordena: la restricción deja 12~clústeres candidatos —las 5~cadenas,
4~\textit{isolated} y 3~limpios—, de modo que las métricas de \textit{ranking} no
se calculan sobre los 86~grafos sino sobre esos 12. Un ordenamiento aleatorio de
ese conjunto obtendría un valor esperado de P@5 de \(5/12 \approx 0{,}42\), que
es la línea de base frente a la que debe leerse el 0,80 obtenido.
\ac{FPR}\_clean\,=\,0,00: ningún clúster limpio aparece
entre los de mayor riesgo. La F1~macro por clasificación \textit{argmax} del
\textit{ensemble} de pliegues sobre el conjunto de test es 0,883 (F1 por clase:
\textit{clean} 1,0; \textit{isolated} 0,981; \textit{attack\_chain} 0,667,
con 3 de las 5~cadenas correctamente clasificadas).

\begin{figure}[htbp]
  \centering
  \caption{Efecto de la restricción de viabilidad estructural}
  \label{fig:structural_gate}
  \includegraphics[width=0.82\textwidth]{fig_structural_gate}

  \fuentepropia
\end{figure}

Dado el tamaño del conjunto de test, estas estimaciones llevan asociados
intervalos de confianza amplios (bootstrap de 10\,000 remuestreos, 95\,\%):
P@5 \([0{,}60,\ 1{,}00]\), \ac{FPR}\_clean \([0{,}00,\ 0{,}40]\) y F1~macro
\([0{,}66,\ 1{,}00]\). Los valores puntuales deben interpretarse como
\textit{consistentes con} los objetivos de diseño, no como demostraciones
concluyentes; la ampliación adicional del corpus de test es la vía directa para
estrechar estos intervalos (TF1, Sección~\ref{sec:trabajo_futuro}). La
Tabla~\ref{tab:ensemble_ablation} recoge las tres métricas de \textit{ranking} del
\textit{ensemble} sobre el conjunto de test.

La Tabla~\ref{tab:ensemble_ablation} desglosa la contribución de cada componente
sobre el conjunto de test (86~grafos, 5~cadenas, techo P@5\,=\,1,00), con la
restricción de viabilidad estructural aplicada
en todas las configuraciones. De la ablación se desprenden tres observaciones:

\begin{enumerate}
  \item La configuración con solo Random Forest degrada gravemente el
    \textit{ranking} (P@1\,=\,0,00, P@5\,=\,0,40). El \textit{risk score} medio de
    los repositorios reales de cadena es \(\approx\)0, de modo que la señal por
    manifiesto no ordena clústeres reales.

  \item La señal de escape por sí sola alcanza P@5\,=\,0,80 gracias a la
    restricción estructural, pero con P@1\,=\,0,00 y P@3\,=\,0,67. Al ser binaria,
    no puede ordenar entre los grafos con capacidad de escape, mayoritarios en la
    región de alto riesgo.

  \item La \ac{GNN} por sí sola iguala tanto a los pesos optimizados por el
    \ac{GA} como a los pesos uniformes: los tres alcanzan 1,00/1,00/0,80. En este
    conjunto de test el \textit{ensemble} no mejora el \textit{ranking} de su
    mejor componente individual, observación que debe leerse con la cautela que
    imponen 5~cadenas de test y que se retoma como limitación en el
    Capítulo~\ref{cap:conclusiones}.
\end{enumerate}

\begin{table}[htbp]
\centering
\caption{Ablación de componentes del \textit{ensemble} en el conjunto de test}
\label{tab:ensemble_ablation}
\begin{tabular}{@{}lrrrr@{}}
\toprule
Configuración & P@1 & P@3 & P@5 & FPR\_clean \\
\midrule
GA-óptimo \((0{,}517/0{,}115/0{,}368)\) & 1,00 & 1,00 & 0,80 & 0,00 \\
Pesos uniformes \((1/3)\)               & 1,00 & 1,00 & 0,80 & 0,00 \\
Solo \ac{GNN}                           & 1,00 & 1,00 & 0,80 & 0,00 \\
Solo Random Forest                            & 0,00 & 0,33 & 0,40 & 0,00 \\
Solo señal de escape                    & 0,00 & 0,67 & 0,80 & 0,00 \\
\bottomrule
\end{tabular}
\fuentepropia
\end{table}

% -----------------------------------------------------------
\section{Análisis de la clasificación por clases: \acs{GNN}}
\label{sec:gnn_por_clase}
% -----------------------------------------------------------

El registro de la validación cruzada conserva por pliegue únicamente la F1 macro,
la Precisión@5 y la exactitud agregada, sin matrices de confusión individuales,
por lo que el desglose por clases se apoya exclusivamente en el conjunto de
test. En él, la matriz de confusión del \textit{ensemble}
de pliegues (Figura~\ref{fig:gnn_confusion}, clasificación \textit{argmax})
concentra el error en la frontera entre cadena y aislado: los dos falsos
negativos de cadena se clasifican como \textit{isolated} y ninguno como
\textit{clean}, de modo que el sistema degrada hacia la clase intermedia y no
hacia la omisión total. Esta observación descansa sobre las cinco cadenas del
conjunto de test y debe leerse con la cautela que impone ese tamaño.

\begin{figure}[htbp]
  \centering
  \caption{Matriz de confusión del \textit{ensemble} de pliegues en test}
  \label{fig:gnn_confusion}
  \includegraphics[width=0.55\textwidth]{fig_gnn_confusion}

  \fuentepropia
\end{figure}

El único fallo de \textit{ranking} del conjunto de test ilustra el error sistemático del
modelo: \texttt{minik8s-ctf} (posición~6, tras el clúster \textit{isolated}
multi-nodo \texttt{kubernetes-extras}) es una cadena de 4~nodos cuya
probabilidad \ac{GNN} es prácticamente nula (0,0003) y cuyo \textit{risk score}
medio del Random Forest es 0,0; su única señal activa es la bandera binaria de escape (la mitad
de sus nodos tienen capacidad de escape). Cuando la
señal estructural que la \ac{GNN} aprende no reconoce el patrón de la cadena y
la señal por manifiesto tampoco lo captura, la bandera de escape binaria por sí
sola no basta para entrar en el top-5. Las cuatro cadenas recuperadas, en
cambio, combinan probabilidad \ac{GNN} alta (0,47--0,63) con fracciones de
escape por nodo muy variables (0,02--0,94; magnitud descriptiva, no la señal
binaria del \textit{ensemble}), lo que confirma que el canal estructural es
el que ordena los repositorios reales. Esta observación motiva el trabajo futuro
TF3 (análisis semántico de \ac{RBAC} completo,
Sección~\ref{sec:trabajo_futuro}).

% -----------------------------------------------------------
\section{Validación con herramientas de análisis estático}
\label{sec:validacion_externa}
% -----------------------------------------------------------

Como contraste externo, se ejecutó Checkov \citep{Checkov2021} sobre el
subconjunto de manifiestos procedentes de GitHub, para el que la herramienta
produjo salida en 1.175 de las 1.510~filas de esa fuente. La descarga previa de
los manifiestos referenciados por Rahman et~al.\ tuvo éxito en 1.985 de 2.039
casos (97,4\,\%); los 54~restantes ya no eran accesibles en origen. El alcance
del contraste queda limitado, por tanto, por la cobertura del escaneo y no por
la disponibilidad de los manifiestos. Para esta validación, el índice \ac{UMI} se calcula como la
proporción de comprobaciones de Checkov que fallan sobre el total
evaluado por manifiesto (\(\in [0,1]\); valores más altos indican mayor densidad
de configuraciones erróneas detectadas). El índice \ac{UMI} medio resultante fue de
0,603, y las comprobaciones incumplidas con mayor frecuencia fueron
\textit{checkov\_no\_run\_as\_nonroot} (38,3\,\%) y
\textit{checkov\_no\_readonly\_rootfs} (37,2\,\%).

El alcance de este contraste es limitado y conviene delimitarlo, porque las dos
fuentes no son independientes. Sobre esas mismas 1.175~filas las características
homólogas del extractor propio (\texttt{NO\_RUN\_AS\_NON\_ROOT} y
\texttt{NO\_READ\_ONLY\_ROOT\_FS}) están ausentes en su totalidad, y es
precisamente la salida de Checkov la que las rellena en tiempo de carga
a través de la tabla de equivalencias \texttt{CHECKOV\_FILL}
(Sección~\ref{sec:construccion_dataset}). Cotejar ambas fuentes manifiesto a
manifiesto equivaldría, por tanto, a comparar esos valores consigo mismos: no
cabe calcular una tasa de acuerdo, y la coincidencia entre las frecuencias
observadas es una identidad de procedencia, no una concordancia entre
detectores. El resultado debe interpretarse, en consecuencia, como una
verificación de cobertura —qué proporción del corpus admite análisis por una
herramienta externa y con qué densidad de incumplimientos— y no como una
validación externa independiente de las predicciones del modelo.

% -----------------------------------------------------------
\section{Métrica global del proceso completo}
\label{sec:metrica_global}
% -----------------------------------------------------------

Las Secciones~\ref{sec:layer1},~\ref{sec:layer2} y~\ref{sec:layer3} evalúan cada
capa por separado, mientras que el comportamiento del sistema de extremo a
extremo —tal como lo experimenta un operador que confía en el \textit{ranking}
final— se resume en una única métrica de proceso completo. Esa
métrica es el conjunto Precisión@k / \ac{FPR}\_clean del sistema desplegado sobre
el conjunto de test, ya presentado en la Sección~\ref{sec:layer3} y recogido en
la Tabla~\ref{tab:ensemble_ablation}: P@1\,=\,1,00, P@3\,=\,1,00, P@5\,=\,0,80 y
\ac{FPR}\_clean\,=\,0,00 sobre 86~grafos con 5~cadenas de ataque de repositorios
reales, no vistos durante el desarrollo; dado el tamaño reducido de
este conjunto, los intervalos de confianza al 95\,\% son amplios (P@5
\(\in\,[0{,}60,\ 1{,}00]\), \ac{FPR}\_clean \(\in\,[0{,}00,\ 0{,}40]\),
Sección~\ref{sec:layer3}), por lo que la cifra debe leerse como indicativa.

Esta cifra es la calidad global de la herramienta porque resume, en un único
indicador operativo, la salida final que recibe el operador tras pasar por las
tres capas: sobre los 86~grafos evaluados, un P@1 de 1,00 indica que la primera
posición del \textit{ranking} correspondió en todos los casos a una cadena de
ataque real, un P@3 de 1,00 extiende esa observación a las tres primeras
posiciones y un \ac{FPR}\_clean nulo indica que ninguno de los tres clústeres
limpios llegó a competir por la atención del operador. Estas cifras se apoyan en
5~cadenas y 3~clústeres limpios, de modo que describen el comportamiento
observado sobre un conjunto reducido y no una garantía del sistema. La
Sección~\ref{sec:layer3} muestra que, sobre este
conjunto de test, el componente \ac{GNN} por sí solo iguala estas cifras y que
el \textit{ensemble} no las mejora; la aportación de la Capa~3 no es superar a
su mejor componente individual, sino producir el mismo resultado sin heredar los
puntos débiles de los demás (la Capa~1 sola colapsa el \textit{ranking} y la señal de
escape sola no puede ordenar por sí misma, Tabla~\ref{tab:ensemble_ablation}).
Esta combinación es lo que hace viable el caso de uso definido en RF-5
(Sección~\ref{sec:requisitos}): priorizar un número reducido de alertas de forma
que la revisión manual de un equipo de seguridad se concentre exclusivamente en
riesgo real.

% -----------------------------------------------------------
\section{Verificación de cumplimiento de requisitos}
\label{sec:verificacion_requisitos}
% -----------------------------------------------------------

El cumplimiento de los requisitos funcionales y no funcionales definidos en el
Capítulo~\ref{cap:requisitos} se verifica a partir de dos clases de evidencia: las
métricas predictivas presentadas en las Secciones~\ref{sec:layer1},
\ref{sec:layer2}, \ref{sec:layer3} y~\ref{sec:metrica_global}, y las
comprobaciones operativas realizadas sobre la herramienta de línea de comandos
descrita en la Sección~\ref{sec:implementacion}. Los requisitos funcionales
RF-1, RF-2 y RF-5 disponen de un criterio de aceptación numérico y se contrastan
contra las métricas ya presentadas, mientras que RF-3, RF-4 y los cuatro
requisitos no funcionales se sustentan en evidencia de ejecución. Ambas clases de
evidencia se recogen en los dos apartados siguientes.

\subsection{Requisitos funcionales}

RF-1 (detección de configuraciones erróneas individuales) queda verificado por los
resultados de la Capa~1 (Sección~\ref{sec:layer1}: F1 macro\,=\,0,9946, cero
falsos negativos en test). RF-2 (detección de cadenas de ataque) queda
verificado con una salvedad: la Capa~2 y el \textit{ensemble} recuperan 4 de las
5~cadenas de ataque del conjunto de test dentro del top-5 del \textit{ranking}
(Secciones~\ref{sec:layer2} y~\ref{sec:layer3}) y clasifican correctamente 3 de
esas 5~cadenas en la matriz de confusión por \textit{argmax}
(F1~\textit{attack\_chain}\,=\,0,667, Sección~\ref{sec:layer3}); el
requisito se cumple para el caso de uso de priorización (RF-5) pero no de forma
perfecta para la clasificación exacta de cada cadena. Ambos requisitos se
confirman además mediante un test de integración de extremo a extremo: dado un
directorio de manifiestos de prueba con banderas de escape conocidas, la salida
del comando \texttt{kubescan scan} produce el veredicto esperado
(\texttt{ATTACK\_CHAIN}) con \textit{ensemble\_score} por encima del umbral de
clasificación; este test forma parte de la suite descrita en la
Sección~\ref{sec:calidad}. RF-5 (priorización de riesgos) queda
verificado por la métrica global de la Sección~\ref{sec:metrica_global}. RF-3 y
RF-4 (interfaz de línea de comandos y formatos de salida) no admiten una
métrica cuantitativa dedicada y se verifican operacionalmente a continuación.

\paragraph{RF-3 — Interfaz de línea de comandos.} El comando
\texttt{kubescan scan <dir>} se ejercitó sobre los casos de uso documentados en el
Anexo~\ref{ap:guia_uso}, con las dos variantes de salida y las opciones que ese
mismo anexo recoge en la Tabla~\ref{tab:cli_options}. El comando \texttt{kubescan live} comparte con
él la ruta de análisis —difiere únicamente en el origen de los manifiestos, que
vuelca del clúster a un directorio temporal— y su invocación queda documentada en
el mismo anexo, si bien no se aporta una transcripción de ejecución sobre un
clúster activo: la verificación de RF-3 se sustenta, por tanto, en la vía de
análisis sobre manifiestos estáticos.

\paragraph{RF-4 — Formatos de salida.} El comando \texttt{kubescan scan} produce
tanto el informe en texto legible como la salida \acs{JSON} estructurada; el
campo \texttt{flags} de esta última, verificado en RNF-1 a continuación, es la
evidencia de que el formato estructurado expone la información necesaria para
su consumo automatizado.

\subsection{Requisitos no funcionales}

Los cuatro requisitos no funcionales definidos en el Capítulo~\ref{cap:requisitos}
se verifican mediante evidencia operativa obtenida sobre la implementación
descrita en el Capítulo~\ref{cap:herramienta}.

\paragraph{Trazabilidad (RNF-1).} El campo \texttt{flags} de la salida \ac{JSON} lista
explícitamente las características de seguridad activas en cada manifiesto, y el
campo \texttt{escape\_capable} indica qué nodos tienen capacidad de escape. Esto
permite al operador de seguridad comprender y auditar las razones concretas de la
clasificación de riesgo producida por el modelo.

\paragraph{Reproducibilidad (RNF-2).} La cadena de resolución de \textit{checkpoints}
descrita en la Sección~\ref{sec:implementacion} garantiza que el mismo directorio de
manifiestos produce siempre el mismo informe dado el mismo conjunto de pesos del
modelo.

\paragraph{Instalabilidad (RNF-3).} El paquete se instaló satisfactoriamente en
entornos Python 3.10 y 3.11 mediante \texttt{pip install -e kubescan/} en sistemas
macOS y Linux (Ubuntu 22.04). Las dependencias principales (PyTorch Geometric,
scikit-learn, Click) se resuelven automáticamente sin compiladores
externos.

\paragraph{Latencia (RNF-4).} Se midió el tiempo de análisis sobre cuatro clústeres
de prueba de diferente tamaño: 12 manifiestos (0,8\,s), 47 manifiestos (1,4\,s),
128 manifiestos (3,1\,s) y 312 manifiestos (7,2\,s) en un equipo con CPU Intel
Core i7 de octava generación sin GPU. El mayor de los clústeres de prueba
disponibles alcanza 312~manifiestos, de modo que el umbral de 500~manifiestos
fijado en el requisito no llegó a medirse de forma directa. La progresión
observada es aproximadamente lineal —en torno a 23\,ms por manifiesto en el punto
de mayor carga—, por lo que la extrapolación a 500~manifiestos sitúa el tiempo
esperado en torno a 12\,s, muy por debajo del límite de 60\,s; el cumplimiento de
RNF-4 se sostiene, por tanto, sobre una extrapolación y no sobre una medición en
el tamaño nominal del requisito.

\section{Evaluación de usabilidad y aplicabilidad con usuarios expertos}
\label{sec:usabilidad}

Además de la verificación de requisitos de la
Sección~\ref{sec:verificacion_requisitos}, se realizó una evaluación de
usabilidad y aplicabilidad de la herramienta con \textbf{usuarios expertos},
siguiendo el protocolo descrito en el repositorio (Anexo~\ref{ap:repositorio}).
Participaron \textbf{tres profesionales} con experiencia en Kubernetes y
\ac{CICD}, a los que se entregó una guía de tareas con órdenes listas para ejecutar y
un cuestionario en línea, en formato autoservicio asíncrono (\(\approx\)10~minutos
por participante). El instrumento se generó de forma reproducible mediante un
\textit{script} versionado (véase Anexo~\ref{ap:repositorio}) y el protocolo se
depuró mediante una prueba piloto simulada, sin participantes reales, cuyo único
objetivo fue verificar la consistencia del cuestionario y del guion de tareas
antes de la sesión con los expertos. La muestra reducida (\(n=3\)) se asume
explícitamente como amenaza a la validez (Sección~\ref{sec:usab_amenazas}); aun
así, tres participantes cubren una fracción sustancial de los problemas de
usabilidad en pruebas formativas \citep{Nielsen2000}.

\subsection{Perfil de los participantes}

La Tabla~\ref{tab:usab_perfil} del Anexo~\ref{ap:repositorio} resume el perfil de
los tres participantes expertos,
seleccionados por su experiencia operando Kubernetes y en seguridad de
contenedores.

\subsection{Éxito por tarea}

Cada participante ejecutó cinco tareas: instalar la herramienta (T1), realizar un
escaneo básico e identificar el veredicto (T2), interpretar el informe con
\texttt{-{}-show-nodes} (T3), producir la salida \ac{JSON} (T4) y, opcionalmente,
escanear un caso propio (T5). La Tabla~\ref{tab:usab_tareas} del
Anexo~\ref{ap:repositorio} recoge la tasa de
éxito (logrado = 1; logrado con ayuda = 0,5; no logrado = 0).

El núcleo del flujo de trabajo —instalación, escaneo e interpretación del
informe— se completó sin ayuda por los tres participantes. Las únicas fricciones
fueron puntuales: un participante necesitó ayuda para extraer campos de la salida
\ac{JSON} (T4) y la tarea opcional sobre un caso propio (T5) se completó con
ayuda en dos casos. La tasa de éxito global del 90\,\%
refleja, por tanto, un flujo de trabajo principal robusto con puntos de fricción
concentrados en tareas secundarias u opcionales.

\subsection{Usabilidad percibida (SUS)}

La usabilidad se midió con la escala \acs{SUS} de diez ítems
\citep{Brooke1996}, que produce una puntuación en el rango
0--100. La media obtenida fue de \textbf{88,3\,\(\pm\)\,10,4} (puntuaciones
individuales de 100, 80 y 85), consistente con la media de referencia del sector
(68) y situada en la banda cualitativa «excelente» de la escala de adjetivos de
\citet{Bangor2009}. Con \(n=3\), sin embargo, el intervalo de confianza al 95\,\%
de la media abarca \([62{,}5;\ 114{,}2]\) y contiene ese valor de referencia, por
lo que los datos no permiten afirmar que la herramienta se sitúe por encima de
él. La puntuación debe interpretarse como indicativa y no como una estimación
precisa del nivel de usabilidad percibida.

\subsection{Aplicabilidad y análisis cualitativo}

Los participantes valoraron cuatro afirmaciones de aplicabilidad (escala 1--5):
integración en \ac{CICD} o auditoría (A1), ajuste de la priorización de riesgo a su
forma de auditar (A2), accionabilidad del informe (A3) y confianza en el
veredicto (A4). Las medias fueron A1\,=\,4,3, A2\,=\,4,3, A4\,=\,4,3 y, como
punto más bajo, A3\,=\,3,7. Con tres respuestas por ítem, estas medias son
sensibles a una única puntuación y conviene acompañarlas de su dispersión: en A3
las puntuaciones individuales fueron 5, 2 y 4, de manera que el descenso de la
media hasta 3,7 procede íntegramente de la valoración de un solo participante
(P2), mientras que los otros dos se sitúan en el rango alto de la escala.
Como fortaleza, un participante destacó exactamente la contribución central del
trabajo: la herramienta «no solo te dice qué está mal en un manifiesto, sino que
conecta varios problemas pequeños para mostrarte si realmente forman un camino de
ataque completo». Ese mismo participante resumió la propuesta, en el apartado de
comentarios finales, como «pensar en cadenas de ataque en lugar de fallos
aislados». Los otros dos destacaron, respectivamente, la interpretación de los
resultados y la rapidez del análisis. La debilidad más recurrente en las respuestas abiertas es
la \textbf{accionabilidad}: el informe comunica el riesgo pero no las acciones
concretas de remediación ni las reglas subyacentes al veredicto. Dos
participantes la mencionaron en el apartado abierto —uno expresó que el
\ac{JSON} «muestra el estado actual pero no las acciones que debería tomar» y
otro la dificultad de «confiar en un puntaje sin conocer bien qué reglas
concretas hay detrás»—, si bien solo uno de ellos la trasladó a la escala
numérica, ya que el segundo puntuó A3 con un 4. La convergencia entre el registro
cualitativo y el cuantitativo es, por tanto, parcial. Esta limitación se recoge como
línea de trabajo futuro (explicabilidad de reglas y recomendaciones de
remediación) en el Capítulo~\ref{cap:conclusiones}. Los escenarios descritos en las
respuestas abiertas fueron dispares: un participante propuso un \textit{gate}
automático en el \textit{pipeline} de \ac{CICD} inmediatamente antes del despliegue
a producción, otro relató haberla aplicado sobre un clúster propio de
preproducción, y el tercero la consideró útil sin concretar un escenario. Solo el
primero de esos usos se corresponde con el paradigma \textit{shift-left} que motiva
el trabajo; la auditoría de clústeres, que el enunciado del ítem~A1 menciona junto
a la integración en \ac{CICD}, no fue propuesta de forma espontánea por ningún
participante.

\subsection{Amenazas a la validez de la evaluación}
\label{sec:usab_amenazas}

El tamaño de la muestra (\(n=3\)) limita la precisión de las estimaciones
puntuales, especialmente de la media SUS, por lo que los resultados se presentan
como indicativos y no concluyentes. Los perfiles cubren desarrollo/IA y
DevOps/SRE pero no un rol de \ac{QA} de \ac{IaC} puro, lo que acota la
generalización. Por último, la participación voluntaria introduce un posible
sesgo de autoselección. Ampliar la muestra y diversificar los perfiles es la vía
directa para robustecer estas conclusiones.
